\section{Casos prácticos y gobernanza de ingeniería}

\subsection{Manipulación robótica de horizonte largo}

El ejemplo de la cocina que presenta Chelsea Finn se compone de varias skill y safety guardrail. Dividimos la tarea en: 1) percepción y estimación de estado, 2) ejecución de la skill (por ejemplo, tomar la olla), 3) verification + recovery.

\begin{knowledgebox}{Puntos de monitoreo del Kitchen Robot}
\begin{itemize}
    \item State check: en cada paso se verifica que gripper/contact funcionen con normalidad;
    \item Skill success flag: al terminar, la skill de bajo nivel reporta `success`;
    \item Recovery template: si la skill falla, el nivel alto activa de inmediato un safe fallback;
    \item Human override: antes de una high-risk action (cutting/pouring) se inserta human-in-the-loop.
\end{itemize}
\end{knowledgebox}

\subsection{Matriz de evaluación y monitoreo}

Se introduce un multi-metric scoreboard para dar seguimiento a accuracy, self-consistency, hallucination rate, token budget, etc. En cada emergent jump hay que revisar varias curvas a la vez para no llegar a un juicio erróneo a partir de una single metric.

\begin{table}[H]
\centering
\begin{tabular}{p{3.5cm}p{5cm}p{3cm}}
\toprule
Métrica & Descripción & Herramienta \\
\midrule
Self-consistency & agreement entre generaciones por múltiples rutas & CoT voting harness \\
Hallucination rate & tasa de errores factuales & human eval + KG check \\
Token budget & consumo de recursos & monitoring hook \\
Instruction drift & desviación de versión del Prompt log & versioned dashboard \\
\bottomrule
\end{tabular}
\caption{Matriz de evaluación del sistema jerárquico.}
\end{table}

\begin{warningbox}{Consecuencias de omitir la cadena de evaluación}
\begin{itemize}
    \item Fijarse solo en accuracy puede ocultar hallucination;
    \item El Prompt drift puede invalidar el high-level plan mientras el nivel bajo sigue ejecutándose;
    \item La falta de monitoreo del token budget provoca compute overrun.
\end{itemize}
\end{warningbox}

\subsection{Incident grade-book y lessons}

\begin{table}[H]
\centering
\begin{tabular}{p{4cm}p{6cm}}
\toprule
Campo & Contenido \\
\midrule
Trigger & hallucination/jump/goal failure \\
Root cause & prompt drift / data shift / skill bug \\
Resolution & rollback prompt / retrain options / add guardrail \\
Next action & drill + doc update \\
\bottomrule
\end{tabular}
\caption{Plantilla del incident grade-book.}
\end{table}

\begin{importantbox}{Procedimiento operativo del Grade-book}
\begin{itemize}
    \item Completar en 24 horas la línea de tiempo trigger → analysis → resolution;
    \item Guardar el snapshot de prompt/data/attention en `notes/incident/`;
    \item Actualizar las lessons en `lessons.md` y compartirlas en la weekly retro.
\end{itemize}
\end{importantbox}

\subsection{Resumen de la sección}

La gobernanza de ingeniería requiere encadenar skill execution, LLM planning y la matriz de monitoreo en un ciclo cerrado, y archivar el prompt log + evaluation snapshot en cada incident.
